\section*{参考文献}
\addcontentsline{toc}{section}{参考文献}
\begin{list}{}{\leftmargin=4.2em\labelwidth=3.8em\labelsep=0.4em\itemindent=0em\parsep=0.35em}
\item[{AB83}] M. F. Atiyah and R. Bott. The Yang-Mills equations over Riemann surfaces. Philos. Trans. Roy. Soc. London Ser. A, 308(1505):523–615, 1983.
\item[{ABC+24a}] D. Arinkin, D. Beraldo, L. Chen, J. Faergeman, D. Gaitsgory, K. Lin, S. Raskin, and N. Rozenblyum. Proof of the geometric Langlands conjecture iv: ambidexterity, 2024. \url{https://arxiv.org/abs/2409.08670}, arXiv: 2409.08670.
\item[{ABC+24b}] Dima Arinkin, Dario Beraldo, Justin Campbell, Lin Chen, Joakim Faergeman, Dennis Gaitsgory, Kevin Lin, Sam Raskin, and Nick Rozenblyum. Proof of the geometric Langlands conjecture ii: Kac-moody localization and the FLE, 2024. \url{https://arxiv.org/abs/2409.07051}, arXiv:2409.07051.
\item[{AF15}] David Ayala and John Francis. Factorization homology of topological manifolds. J. of Topology, 8(4):1045–1084, 2015. doi:10.1112/jtopol/jtv028.
\item[{AG15}] Dima Arinkin and Dennis Gaitsgory. Singular support of coherent sheaves, and the geometric Langlands conjecture. Selecta Math. (N.S.), 21(1):1–199, 2015. doi:10.1007/s00029-014-0167-5.
\item[{AGK+20}] Dima Arinkin, Dennis Gaitsgory, David Kazhdan, Sam Raskin, Nick Rozenblyum, and Yakov Varshavsky. The stack of local systems with restricted variation and geometric Langlands theory with nilpotent singular support. 2020. arXiv:2010.01906.
\item[{AGK+21}] Dima Arinkin, Dennis Gaitsgory, David Kazhdan, Sam Raskin, Nick Rozenblyum, and Yakov Varshavsky. Automorphic functions as the trace of Frobenius. 2021. arXiv:2102.07906.
\item[{AR18}] Pramod N. Achar and Simon Riche. Reductive groups, the loop Grassmannian, and the Springer resolution. Invent. Math., 214(1):289–436, 2018.
\item[{Ari05}] D. Arinkin. On λ-connections on a curve where λ is a formal parameter. Math. Res. Lett., 12(4):551–565, 2005.
\item[{Ari13}] Dima Arinkin. Autoduality of compactified Jacobians for curves with plane singularities. J. Algebraic Geom., 22(2):363–388, 2013. \url{https://doi.org/10.1090/S1056-3911-2012-00596-7}, doi:10.1090/S1056-3911-2012-00596-7.
\item[{Ari16}] Dima Arinkin. Irreducible connections admit generic oper structures, 2016. \url{https://arxiv.org/abs/1602.08989}, arXiv:1602.08989.
\item[{BB07}] Alexander Braverman and Roman Bezrukavnikov. Geometric Langlands correspondence for D-modules in prime characteristic: the $GL\_n(q)$ case. Pure Appl. Math. Q., 3(1):153–179, 2007.
\item[{BD}] Alexander Beilinson and Vladimir Drinfeld. Quantization of Hitchin Hamiltonians and Hecke eigensheaves. \url{https://math.uchicago.edu/~drinfeld/Langlands/QuantizationHitchin.pdf.}
\item[{BD04}] Alexander Beilinson and Vladimir Drinfeld. Chiral algebras, volume 51 of American Mathematical Society Colloquium Publications. American Mathematical Society, Providence, RI, 2004. doi:10.1090/coll/051.
\item[{Bez06}] Roman Bezrukavnikov. Noncommutative counterparts of the Springer resolution. In International Congress of Mathematicians. Vol. II, pages 1119–1144. Eur. Math. Soc., Zürich, 2006.
\item[{Bez16}] Roman Bezrukavnikov. On two geometric realizations of an affine Hecke algebra. Publ. Math. Inst. Hautes Études Sci., 123:1–67, 2016. arXiv:1209.0403.
\item[{BF08}] Roman Bezrukavnikov and Mikhail Finkelberg. Equivariant Satake category and Kostant-Whittaker reduction. Mosc. Math. J. 8 (2008), no. 1, 39–72, 183., 8(1):39–72, 2008. doi:10.17323/1609-4514-2008-8-1-39-72.
\item[{BKS22}] Dario Beraldo, David Kazhdan, and Tomer M. Schlank. Contractibility of the space of generic opers for classical groups, 2022. \url{https://arxiv.org/abs/1801.00655}, arXiv:1801.00655.
\item[{BM13}] Roman Bezrukavnikov and Ivan Mirković. Representations of semisimple Lie algebras in prime characteristic and the noncommutative Springer resolution. Ann. of Math. (2), 178(3):835–919, 2013. With an appendix by Eric Sommers.
\item[{BMR08}] Roman Bezrukavnikov, Ivan Mirković, and Dmitriy Rumynin. Localization of modules for a semisimple Lie algebra in prime characteristic. Ann. of Math. (2), 167(3):945–991, 2008. With an appendix by Bezrukavnikov and Simon Riche.
\item[{BZN18}] David Ben-Zvi and David Nadler. Betti geometric Langlands. Proc. Sympos. Pure Math., 97.2, 2018. doi: 10.1090/pspum/097.2/01.
\item[{BZSV24}] David Ben-Zvi, Yiannis Sakellaridis, and Akshay Venkatesh. Relative Langlands duality, 2024. \url{https://arxiv.org/abs/2409.04677}, arXiv:2409.04677.
\item[{Car}] Ana Caraiani. An excursion through the land of shtukas. \url{https://www.ams.org/meetings/lectures/current-events-bulletin.}
\item[{Car06}] R. W. Carter. A survey of the work of George Lusztig. Nagoya Math. J., 182:1–45, 2006.
\item[{CCF+24}] Justin Campbell, Lin Chen, Joakim Faergeman, Dennis Gaitsgory, Kevin Lin, Sam Raskin, and Nick Rozenblyum. Proof of the geometric Langlands conjecture iii: compatibility with parabolic induction, 2024. \url{https://arxiv.org/abs/2409.07051}, arXiv:2409.07051.
\item[{CDR21}] Justin Campbell, Gurbir Dhillon, and Sam Raskin. Fundamental local equivalences in quantum geometric Langlands. Compos. Math., 157(12):2699–2732, 2021.
\item[{CF21}] Lin Chen and Yuchen Fu. An extension of the Kazhdan-Lusztig equivalence, 2021. \url{https://arxiv.org/abs/2111.14606}, arXiv:2111.14606.
\item[{CG17}] Kevin Costello and Owen Gwilliam. Factorization algebras in quantum field theory. Vol. 1, volume 31 of New Mathematical Monographs. Cambridge University Press, Cambridge, 2017. doi:10.1017/9781316678626.
\item[{CGKS25}] Pierre-Henri Chaudouard, Wee Gan, Tasho Kaletha, and Yiannis Sakellaridis, editors. The Langlands Program, volume 112.2 of Proceedings of Symposia in Pure Mathematics. American Mathematical Society, Providence, RI, 2025.
\item[{Cor88}] Kevin Corlette. Flat G-bundles with canonical metrics. J. Differential Geom., 28(3):361–382, 1988. \url{http://projecteuclid.org/euclid.jdg/1214442469.}
\item[{CZ17}] Tsao-Hsien Chen and Xinwen Zhu. Geometric Langlands in prime characteristic. Compos. Math., 153(2):395– 452, 2017.
\item[{DG02}] Ron Y. Donagi and Dennis Gaitsgory. The gerbe of Higgs bundles. Transform. Groups, 7(2):109–153, 2002.
\item[{DG13}] Vladimir Drinfeld and Dennis Gaitsgory. On some finiteness questions for algebraic stacks. Geom. Funct. Anal, pages 149–294, 2013. doi:10.1007/s00039-012-0204-5.
\item[{DG15}] Vladimir Drinfeld and Dennis Gaitsgory. Compact generation of the category of D-modules on the stack of G-bundles on a curve. Camb. J. Math., 3(1-2):19–125, 2015. doi:10.4310/CJM.2015.v3.n1.a2.
\item[{DM96}] Ron Donagi and Eyal Markman. Spectral covers, algebraically completely integrable, Hamiltonian systems, and moduli of bundles. In Integrable systems and quantum groups (Montecatini Terme, 1993), volume 1620 of Lecture Notes in Math., pages 1–119. Springer, Berlin, 1996.
\item[{DP09}] R. Donagi and T. Pantev. Geometric Langlands and non-abelian Hodge theory. In Surveys in differential geometry. Vol. XIII. Geometry, analysis, and algebraic geometry: forty years of the Journal of Differential Geometry, volume 13 of Surv. Differ. Geom., pages 85–116. Int. Press, Somerville, MA, 2009.
\item[{DP12}] Ron Donagi and Tony Pantev. Langlands duality for Hitchin systems. Invent. Math., 189(3):653–735, 2012.
\item[{DP22}] Ron Donagi and Tony Pantev. Parabolic Hecke eigensheaves. Astérisque, (435):viii+192, 2022.
\item[{DPS24}] Ron Donagi, Tony Pantev, and Carlos Simpson. Twistor hecke eigensheaves in genus 2, 2024. \url{https://arxiv.org/abs/2403.17045}, arXiv:2403.17045.
\item[{Dri83}] V. G. Drinfeld. Two-dimensional l-adic representations of the fundamental group of a curve over a finite field and automorphic forms on $\mathrm{GL}_{2}(\F_{q})$. Amer. J. Math., 105(1):85–114, 1983.
\item[{EFK21}] Pavel Etingof, Edward Frenkel, and David Kazhdan. An analytic version of the Langlands correspondence for complex curves. In Integrability, quantization, and geometry II. Quantum theories and algebraic geometry, volume 103.2 of Proc. Sympos. Pure Math., pages 137–202. Amer. Math. Soc., Providence, RI, [2021] ©2021. \url{https://doi.org/10.1090/pspum/103.2/01856}, doi:10.1090/pspum/103.2/01856. 32
\item[{FF92}] Boris Feigin and Edward Frenkel. Affine Kac-Moody algebras at the critical level and Gelfand-Dikii algebras. In Infinite analysis, Part A, B (Kyoto, 1991), volume 16 of Adv. Ser. Math. Phys., pages 197–215. World Sci. Publ., River Edge, NJ, 1992.
\item[{FG06}] Vladimir Fock and Alexander Goncharov. Moduli spaces of local systems and higher Teichmüller theory. Publ. Math. Inst. Hautes Études Sci., (103):1–211, 2006.
\item[{FG09a}] Edward Frenkel and Dennis Gaitsgory. Local geometric Langlands correspondence: the spherical case. In Algebraic analysis and around, volume 54 of Adv. Stud. Pure Math., pages 167–186. Math. Soc. Japan, Tokyo, 2009. \url{https://doi.org/10.2969/aspm/05410167}, doi:10.2969/aspm/05410167.
\item[{FG09b}] Edward Frenkel and Dennis Gaitsgory. Localization of g-modules on the affine Grassmannian. Ann. of Math. (2), 170(3):1339–1381, 2009.
\item[{FG09c}] Edward Frenkel and Benedict Gross. A rigid irregular connection on the projective line. Ann. of Math. (2), 170(3):1469–1512, 2009. doi:10.4007/annals. 2009.170.1469.
\item[{FG20}] Edward Frenkel and Davide Gaiotto. Quantum Langlands dualities of boundary conditions, D-modules, and conformal blocks. Commun. Number Theory Phys., 14(2):199–313, 2020. doi:10.4310/CNTP.2020.v14.n2.a1.
\item[{FGV01}] Edward Frenkel, Dennis Gaitsgory, and Kari Vilonen. Whittaker patterns in the geometry of moduli spaces of bundles on curves. Ann. Math, pages 699–748, 2001. doi:10.2307/2661366.
\item[{Fre04}] Edward Frenkel. Recent advances in the Langlands program. Bull. Amer. Math. Soc. (N.S.), 41(2):151–184, 2004. doi:10.1090/S0273-0979-04-01001-8.
\item[{Fre05}] Edward Frenkel. Gaudin model and opers. In Infinite dimensional algebras and quantum integrable systems, volume 237 of Progr. Math., pages 1–58. Birkhäuser, Basel, 2005. doi:10.1007/3-7643-7341-5\_1.
\item[{Fre07}] Edward Frenkel. Langlands correspondence for loop groups, volume 103 of Cambridge Studies in Advanced Mathematics. Cambridge University Press, Cambridge, 2007.
\item[{Fre10}] Edward Frenkel. Gauge theory and Langlands duality. Astérisque, (332):Exp. No. 1010, ix–x, 369–403, 2010. Séminaire Bourbaki. Volume 2008/2009. Exposés 997–1011.
\item[{FrR25}] Joakim Færgeman and Sam Raskin. Non-vanishing of geometric Whittaker coefficients for reductive groups. J. Amer. Math. Soc., 38(4):919–995, 2025.
\item[{FS21}] Laurent Fargues and Peter Scholze. Geometrization of the local Langlands correspondence. 2021. arXiv:2102.13459.
\item[{Gai04}] D. Gaitsgory. On a vanishing conjecture appearing in the geometric Langlands correspondence. Ann. of Math. (2), 160(2):617–682, 2004. \url{https://doi.org/10.4007/annals. 2004.160.617}, doi:10.4007/annals. 2004.160.617.
\item[{Gai07}] D. Gaitsgory. On de Jong’s conjecture. Israel J. Math., 157:155–191, 2007.
\item[{Gai08}] D. Gaitsgory. Twisted Whittaker model and factorizable sheaves. Selecta Math. (N.S.), 13(4):617–659, 2008.
\item[{Gai15}] Dennis Gaitsgory. Outline of the proof of the geometric Langlands conjecture for GL2. Astérisque, (370):1–112, 2015.
\item[{Gai16}] Dennis Gaitsgory. Functors given by kernels, adjunctions and duality. J. Algebraic Geom., 25(3):461–548, 2016. doi:10.1090/jag/654.
\item[{Gai17a}] Dennis Gaitsgory. Progrès récents dans la théorie de Langlands géométrique. Number 390, pages Exp. No. 1109, 139–168, 2017. Séminaire Bourbaki. Vol. 2015/2016. Exposés 1104–1119.
\item[{Gai17b}] Dennis Gaitsgory. A “strange” functional equation for Eisenstein series and miraculous duality on the moduli stack of bundles. Ann. Sci. Éc. Norm. Supér. (4), 50(5):1123–1162, 2017. doi:10.24033/asens. 2341.
\item[{Gai25}] Dennis Gaitsgory. Local and global Langlands conjecture(s) over function fields. to appear, proceedings of the 2026 icm, 2025. \url{https://arxiv.org/abs/2509.24902}, arXiv:2509.24902.
\item[{Gan}] Wee Teck Gan. Relative Langlands Duality, after Ben-Zvi, Sakellaridis and Venkatesh. Seminaire Bourbaki 29 Mars 2025, \url{https://www.bourbaki.fr/seminaires/2024-25/Prog\_mar-25.html.}
\item[{Gin}] Victor Ginzburg. Perverse sheaves on a loop group and Langlands’ duality. arXiv:alg-geom/9511007v4.
\item[{GL19a}] D. Gaitsgory and S. Lysenko. Metaplectic Whittaker category and quantum groups: the “small” FLE, 2019. \url{https://arxiv.org/abs/1903.02279}, arXiv:1903.02279.
\item[{GL19b}] Dennis Gaitsgory and Jacob Lurie. Weil’s Conjecture for Function Fields: Volume I (AMS-199). Princeton University Press, 2019.
\item[{GMN13}] Davide Gaiotto, Gregory W. Moore, and Andrew Neitzke. Wall-crossing, Hitchin systems, and the WKB approximation. Adv. Math., 234:239–403, 2013.
\item[{GPS24}] Sheel Ganatra, John Pardon, and Vivek Shende. Microlocal Morse theory of wrapped Fukaya categories. Ann. of Math. (2), 199(3):943–1042, 2024.
\item[{GR17}] Dennis Gaitsgory and Nick Rozenblyum. A study in derived algebraic geometry. Vol. I. Correspondences and duality, Vol. II. Deformations, Lie theory and formal geometry, volume 221 of Mathematical Surveys and Monographs. American Mathematical Society, Providence, RI, 2017.
\item[{GR24a}] Dennis Gaitsgory and Sam Raskin. Proof of the geometric Langlands conjecture i: construction of the functor, 2024. \url{https://arxiv.org/abs/2405.03599}, arXiv:2405.03599.
\item[{GR24b}] Dennis Gaitsgory and Sam Raskin. Proof of the geometric Langlands conjecture v: the multiplicity one theorem, 2024. \url{https://arxiv.org/abs/2409.09856}, arXiv:2409.09856.
\item[{GR25}] Dennis Gaitsgory and Sam Raskin. Geometric Langlands in positive characteristic from characteristic zero, 2025. \url{https://arxiv.org/abs/2508.02237}, arXiv:2508.02237.
\item[{GW09a}] Davide Gaiotto and Edward Witten. S-duality of boundary conditions in N “ 4 super Yang-Mills theory. Adv. Theor. Math. Phys., 13(3):721–896, 2009.
\item[{GW09b}] Davide Gaiotto and Edward Witten. Supersymmetric boundary conditions in N “ 4 super Yang-Mills theory. J. Stat. Phys., 135(5-6):789–855, 2009. doi:10.1007/s10955-009-9687-3.
\item[{GW24}] Davide Gaiotto and Edward Witten. Gauge theory and the analytic form of the geometric Langlands program. Ann. Henri Poincaré, 25(1):557–671, 2024. \url{https://doi.org/10.1007/} s00023-022-01225-6, doi:10.1007/s00023-022-01225-6.
\item[{Hal12}] Thomas C. Hales. The fundamental lemma and the Hitchin fibration [after Ngˆo Bao Chˆau]. Number 348, pages Exp. No. 1035, ix, 233–263.2012. Séminaire Bourbaki: Vol. 2010/2011. Exposés 1027–1042.
\item[{Hau23}] Tamás Hausel. Enhanced mirror symmetry for Langlands dual Hitchin systems. In ICM—International Congress of Mathematicians. Vol. 3. Sections 1–4, pages 2228–2249. EMS Press, Berlin, [2023] ©2023.
\item[{Hit87a}] N. J. Hitchin. The self-duality equations on a Riemann surface. Proc. London Math. Soc. (3), 55(1):59–126, 1987.
\item[{Hit87b}] Nigel Hitchin. Stable bundles and integrable systems. Duke Math. J., 54(1):91–114, 1987.
\item[{Hit92}] N. J. Hitchin. Lie groups and Teichmüller space. Topology, 31(3):449–473, 1992.
\item[{HNY13}] Jochen Heinloth, Bao-Chˆau Ngˆo, and Zhiwei Yun. Kloosterman sheaves for reductive groups. Ann. of Math. (2), 177(1):241–310, 2013. doi:10.4007/annals. 2013.177.1.5.
\item[{HT03}] Tamás Hausel and Michael Thaddeus. Mirror symmetry, Langlands duality, and the Hitchin system. Invent. Math., 153(1):197–229, 2003. \url{https://doi.org/10.1007/s00222-003-0286-7}, doi:10.1007/s00222-003-0286-7.
\item[{Jor23}] David Jordan. Langlands duality for skein modules of 3-manifolds, 2023. \url{https://arxiv.org/abs/2302.14734}, arXiv:2302.14734.
\item[{KL93}] D. Kazhdan and G. Lusztig. Tensor structures arising from affine Lie algebras. I, II. J. Amer. Math. Soc., 6(4), 1993.
\item[{Kon09}] Maxim Kontsevich. Symplectic geometry of homological algebra. 2009. \url{https://archive.mpim-bonn.mpg.} de/id/eprint/1536/1/preprint\_2009\_40\_a.pdf.
\item[{KW07}] Anton Kapustin and Edward Witten. Electric-magnetic duality and the geometric Langlands program. Commun. Num. Theor. Phys., 1:1–236, 2007. doi:10.4310/CNTP.2007.v1.n1.a1.
\item[{Laf18}] Vincent Lafforgue. Shtukas for reductive groups and Langlands correspondence for function fields. In Proceedings of the International Congress of Mathematicians—Rio de Janeiro 2018. Vol. I. Plenary lectures, pages 635–668. World Sci. Publ., Hackensack, NJ, 2018.
\item[{Lau87}] Gérard Laumon. Correspondance de Langlands géométrique pour les corps de fonctions. Duke Math. J., 54(2):309–359, 1987.
\item[{Lau88}] Gérard Laumon. Un analogue global du cˆone nilpotent. Duke Math. J., 57(2):647–671, 1988.
\item[{Lau96}] Gerard Laumon. Transformation de Fourier generalisée, 1996. arXiv:alg-geom/9603004.
\item[{LNY24}] Penghui Li, David Nadler, and Zhiwei Yun. Functions on the commuting stack via Langlands duality. Ann. of Math. (2), 200(2):609–748, 2024. doi:10.4007/annals. 2024.200.2.5.
\item[{LT24}] Thomas Lam and Nicolas Templier. The mirror conjecture for minuscule flag varieties. Duke Math. J., 173(1):75– 175, 2024. doi:10.1215/00127094-2024-0007.
\item[{Lur}] Jacob Lurie. Higher algebra. \url{http://math.ias.edu/~lurie/papers/HA.pdf.}
\item[{Lur09}] Jacob Lurie. Higher topos theory, volume 170 of Annals of Mathematics Studies. Princeton University Press, Princeton, NJ, 2009. \url{http://math.ias.edu/~lurie/papers/HTT.pdf.}
\item[{Lus83}] George Lusztig. Singularities, character formulas, and a q-analog of weight multiplicities. 101:208–229, 1983.
\item[{Mor12}] Masanori Morishita. Knots and primes. Universitext. Springer, London, 2012. An introduction to arithmetic topology.
\item[{MRV19a}] Margarida Melo, Antonio Rapagnetta, and Filippo Viviani. Fourier-Mukai and autoduality for compactified Jacobians. I. J. Reine Angew. Math., 755:1–65, 2019.
\item[{MRV19b}] Margarida Melo, Antonio Rapagnetta, and Filippo Viviani. Fourier-Mukai and autoduality for compactified Jacobians, II. Geom. Topol., 23(5):2335–2395, 2019.
\item[{Muk81}] Shigeru Mukai. Duality between $D(X)$ and $D(\hat{X})$ with its application to Picard sheaves. Nagoya Math. J., 81:153–175, 1981. \url{http://projecteuclid.org/euclid.nmj/1118786312.}
\item[{MV07}] Ivan Mirković and Kari Vilonen. Geometric Langlands duality and representations of algebraic groups over commutative rings. Ann. of Math. (2), 166(1):95–143, 2007. doi:10.4007/annals. 2007.166.95.
\item[{Nad09}] David Nadler. Microlocal branes are constructible sheaves. Selecta Math. (N.S.), 15(4):563–619, 2009.
\item[{Nad12}] David Nadler. The geometric nature of the fundamental lemma. Bull. Amer. Math. Soc. (N.S.), 49(1):1–50, 2012. doi:10.1090/S0273-0979-2011-01342-8. 34
\item[{NT25}] David Nadler and Jeremy Taylor. The Whittaker Functional Is a Shifted Microstalk. Transform. Groups, 30(3):1425–1450, 2025.
\item[{NY19}] David Nadler and Zhiwei Yun. Spectral action in Betti geometric Langlands. Israel J. Math., 232(1):299–349, 2019. doi:10.1007/s11856-019-1871-9.
\item[{NZ09}] David Nadler and Eric Zaslow. Constructible sheaves and the Fukaya category. J. Amer. Math. Soc., 22(1):233– 286, 2009. doi:10.1090/S0894-0347-08-00612-7.
\item[{PT25}] Tudor Pădurariu and Yukinobu Toda. The dolbeault geometric Langlands conjecture via limit categories, 2025. \url{https://arxiv.org/abs/2508.19624}, arXiv:2508.19624.
\item[{Ras}] Sam Raskin. Conjectures of arthur and ramanujan for unramified automorphic forms: Announcement. to appear, proceedings of the 2026 icm.
\item[{Ric24}] Simon Riche. Some applications of the geometric satake equivalence to modular representation theory, 2024. \url{https://arxiv.org/abs/2403.03734}, arXiv:2403.03734.
\item[{Rot96}] Mitchell Rothstein. Sheaves with connection on abelian varieties. Duke Math. J., 84(3):565–598, 1996. doi: 10.1215/S0012-7094-96-08418-5.
\item[{RW22}] Simon Riche and Geordie Williamson. Smith-Treumann theory and the linkage principle. Publ. Math. Inst. Hautes Études Sci., 136:225–292, 2022.
\item[{RY24}] Sam Raskin and David Yang. Affine Beilinson-Bernstein localization at the critical level. Ann. of Math. (2), 200(2):487–527, 2024.
\item[{Sch}] Peter Scholze. Geometrization of the real local Langlands correspondence. \url{https://people.mpim-bonn.mpg.de/} scholze/RealLocalLanglands.pdf.
\item[{Sch25}] Peter Scholze. Geometrization of the local Langlands correspondence, motivically, 2025. \url{https://arxiv.} org/abs/2501.07944, arXiv:2501.07944.
\item[{Sch26}] Peter Scholze. Geometric Langlands (after gaitsgory, raskin, ...).2026. \url{https://people.mpim-bonn.mpg.de/} scholze/Exp1252\_Scholze.pdf.
\item[{Sim92}] Carlos T. Simpson. Higgs bundles and local systems. Inst. Hautes Études Sci. Publ. Math., (75):5–95, 1992. \url{http://www.numdam.org/item?id=PMIHES\_1992\_\_75\_\_5\_0.}
\item[{Str17}] Benoˆıt Stroh. La paramétrisation de Langlands globale sur les corps de fonctions. Number 390, pages Exp. No. 1110, 169–197.2017. Séminaire Bourbaki. Vol. 2015/2016. Exposés 1104–1119.
\item[{SV17}] Yiannis Sakellaridis and Akshay Venkatesh. Periods and harmonic analysis on spherical varieties. Astérisque, (396):360, 2017.
\item[{Wil17}] Geordie Williamson. Algebraic representations and constructible sheaves. Jpn. J. Math., 12(2):211–259, 2017.
\item[{XZ22}] Daxin Xu and Xinwen Zhu. Bessel F-isocrystals for reductive groups. Invent. Math., 227(3):997–1092, 2022. doi:10.1007/s00222-021-01079-5.
\item[{Yun14}] Zhiwei Yun. Motives with exceptional Galois groups and the inverse Galois problem. Invent. Math., 196(2):267– 337, 2014. doi:10.1007/s00222-013-0469-9.
\item[{Yun17}] Zhiwei Yun. Rigidity in the Langlands correspondence and applications. In Proceedings of the Sixth International Congress of Chinese Mathematicians. Vol. I, volume 36 of Adv. Lect. Math. (ALM), pages 199–234. Int. Press, Somerville, MA, 2017.
\item[{Zhu23}] Xinwen Zhu. Arithmetic and geometric Langlands program. In ICM—International Congress of Mathematicians. Vol. 3. Sections 1–4, pages 2012–2045. EMS Press, Berlin, [2023] ©2023. University of Texas at Austin Email address: benzvi@utexas.edu
\end{list}